\section{The wrong direction}
\label{sec:lvillposed}

Formula \eqref{eq:lvextract} invites a procedure: interpolate the quoted
smile, differentiate it twice in strike and once in maturity, divide.  On a
continuum of exact prices this is a theorem.  On a real chain --- a few
dozen noisy quotes per expiry --- it fails, and the failure is structural,
not bad luck.

Quantify it.  Suppose each quoted total variance carries an error of size
$\varepsilon$, and strikes are spaced $\Delta k$ apart.  A second difference
estimates curvature as
\[
\partial_{kk}w\;\approx\;
\frac{w(k-\Delta k)-2\,w(k)+w(k+\Delta k)}{\Delta k^{2}},
\]
and the worst configuration of errors --- alternating signs, which is
exactly what a bid--ask bounce produces --- enters this estimate at size
$4\varepsilon/\Delta k^{2}$.  The noise is amplified by the \emph{inverse square}
of the strike spacing: a finer quote grid makes the estimate worse, without
bound.  There is no stable limit to converge to; numerical differentiation
of noisy data is the canonical ill-posed operation.  And the place this
amplified noise lands could not be worse chosen: by \eqref{eq:lvratio} the
curvature sits inside $g_{\rm D}$, the sign-carrying factor of the
risk-neutral density.  The first
time the noise pushes the estimated density through zero, the extraction
returns a \emph{negative variance} --- not a large error but a
non-object, which no pricing equation, no simulation and no risk system
downstream can absorb.

\begin{figure}[htbp]
\centering
\paperfig{\textwidth}{fig_lv_wrongway.pdf}
\caption{The same synthetic market, read both ways.  Quotes are generated
from a known smooth field (dashed) and carry a deterministic alternating
$\pm$\MacLvWrongRippleBp\,bp volatility ripple --- a stylized bid--ask
bounce.  (a)~The extraction \eqref{eq:lvratio} from the noisy quotes at
$\tau=0.25$.  From 13 quotes per expiry it saw\-tooths between
\MacLvWrongCoarseMinPct\% and \MacLvWrongMaxSpikePct\% around a true ATM
level of \MacLvWrongTruthAtmPct\%; from 31 quotes of the \emph{same}
market --- finer $\Delta k$, same $\pm$\MacLvWrongRippleBp\,bp ---
it fails outright at \MacLvWrongNBad\ strikes (crosses: negative variance
or negative density), exactly the $4\varepsilon/\Delta k^{2}$ scaling at work.
(b)~The same 31-quote noisy chains handed to the forward calibration of
this chapter: an admissible surface by construction, within
\MacLvWrongFwdMaxPts\ volatility points of the truth along the plotted
cross-sections, with reprice residuals of rms
\MacLvWrongQuoteRmsBp\,bp --- the scale of the ripple itself.  The noise
stays in the residual, where it belongs.}
\label{fig:lvwrongway}
\end{figure}

\Cref{fig:lvwrongway} performs the experiment on a synthetic market whose
answer is known (the field and quote recipe are in \cref{app:lvnumerics}).
The left panel is the naive pipeline at two quote densities and needs
little commentary beyond its caption: refining the strike grid, which
improves every well-posed computation, makes this one \emph{worse},
quadratically, until half the strikes report that no diffusion exists.
The right panel previews the alternative developed in the rest of the
chapter.  The same noisy quotes are handed to a calibration that never
differentiates them; every iterate it produces is a genuine positive
field; and the fitted surface tracks the truth through the quoted region
while its residuals absorb the ripple.  One caveat, quantified now and
explained in \cref{sec:lvinfo}: at the edges of the shortest expiry's
quoted span, where a single rippled quote is the only witness, the fitted
field wanders by as much as \MacLvWrongEdgePts\ volatility points.  Noise
rejection is a property of dense data and declared smoothing.

The repair is a change of role, not of formula, and it is due to Dupire's
equation being read in its \emph{forward} direction, as Andreasen and Huge
taught \citep{AndreasenHuge2011}.  Decide that the unknown is the field
$v$ itself.  Parametrize it.  Price it through \eqref{eq:lvdupire} ---
which differentiates only the smooth numerical solution, never the data
--- and ask an optimizer for the positive field whose prices meet the
quotes.  Positivity is imposed on the \emph{parameters}, so every iterate
of the search is an admissible diffusion before any quote is consulted.
Readers of Chapters~2--3 will recognize the design: it is the same
constraint-elimination idea that built the LQD chart and the structural
SVI chart (\cref{sec:charts,sec:svistructural}), one dimension up ---
choose the coordinates in which the constraint set is trivial, and let the
pricing map do the hard work.  Three questions remain, and they organize
the rest of the chapter: what parametrization makes positivity trivial in
two dimensions (\cref{sec:lvsheet}); what discretization of
\eqref{eq:lvdupire} keeps \cref{prop:lvpropagate} alive on a finite grid
(\cref{sec:lvlattice}); and what, exactly, the quotes determine about the
field once the formula \eqref{eq:lvextract} is off the table
(\cref{sec:lvinfo}).
